\section{Conclusion}\label{sec:conclusion}

A conformal soft sensor recalibrated with late laboratory labels is an integral loop with dead
time. Its linearisation bounds the step by $2\sin(\pi/(4\theta+2))$, its marginal coverage
converges however many of its intervals are uninformative, and on persistent process residuals
the step that holds a budget lies far below the bound. In a pre-registered study on the Tennessee
Eastman process, a step tuned without the delay left 6.9\,\% of the intervals uninformative at
nominal coverage. Replaying the calibration residuals through the delayed loop held a 2\,\%
budget, and compensating the residual over the dead time narrowed the intervals by 18\,\% at the
same budget, made the misses less persistent and admitted a larger step; no baseline
configuration was both valid and narrower. The width ratio that compensation buys was predicted
from the calibration residuals' autocorrelation to within 5\,\% on average at four updates, and
the gain it implies was never overstated on average, so an engineer can price a faster or closer
analyzer before buying it. The diagnosis
did not appear on the debutanizer stream, where split conformal covered only 0.750 and the
compensated loop 0.900. Random and irregular delays, plant data with time-stamped laboratory
delays, and a step-size-free loop analysed under delay are the next steps.
